\documentclass[aps,prd,reprint,nofootinbib]{revtex4-2}
\usepackage[T1]{fontenc}
\usepackage{amsmath,amssymb,bm}
\usepackage[colorlinks=true,linkcolor=blue,citecolor=blue,urlcolor=blue]{hyperref}

\newcommand{\dd}{\mathrm{d}}
\newcommand{\ve}{\varepsilon}
\newcommand{\reff}{\rho_{\mathrm{eff}}}
\newcommand{\Meff}{M_{\mathrm{eff}}}

\begin{document}

\title{Lattice depletion, signal time advances, and effective repulsion}
% Author and affiliation deliberately omitted pending author instructions.
\date{October 3, 2026}

\begin{abstract}
We consider a lattice modification in which intermediate processing nodes
are removed and the remaining nodes are reconnected, so that a signal
crossing a region encounters fewer elementary delays than in a reference
lattice. Under stated continuum and isotropy assumptions, the relative
encounter count becomes an optical index $n$. Depletion, $n<1$, produces
a travel-time advance; a radially monotone central depletion deflects
passing rays outward. These results concern optical geometry and do not require
assigning negative energy to missing constituents. A gravitational
interpretation additionally requires a prescription for clocks and
massive probes. With the usual static weak-field metric, the same
depletion corresponds to a positive potential and repulsive free fall.
A smooth example makes the distinction explicit, while a negative-mass
type exterior requires a $1/r$ encounter deficit. The model formalizes
the intuition that fewer processing steps reverse the optical effect
of additional delays, and identifies the further assumptions needed to
interpret that reversal as gravity.
\end{abstract}

\maketitle

\section{The proposed operation}
\label{sec:operation}

Suppose that signal propagation through a microscopic substrate takes
time because a signal undergoes successive local processing events.
Adding events along a route then increases its traversal time. The
question considered here is the converse: what would observers built
from the substrate experience if a traversable route required fewer
events?

The intended operation is \emph{depletion with reconnection}. A sequence
of processing nodes is shortened while passage between its surviving
endpoints is retained. The delay associated with a removed node is
omitted; it is not reassigned to the new connection. Thus a region with
fewer nodes can be crossed earlier than the corresponding region of an
undepleted reference lattice. This propagation rule is the defining
physical assumption of the model.

Geometric descriptions of lattice defects provide a broader setting
for this question~\cite{kleinert-2004}. For embedded observers we follow
Toffoli's discussion of observers within uniform systems~\cite{toffoli:79}.
Here the observer and its measuring devices belong to the same
substrate. The processing rule specified here is an additional
assumption. We derive its optical
consequences before introducing a gravitational interpretation. The
distinction between a propagation metric and gravitational dynamics is
also familiar from acoustic analogues~\cite{VisserAnalog}.

The reference lattice supplies the comparison meant by ``fewer.'' Its
coordinates label corresponding regions before and after modification;
they need not coincide with the distances inferred inside the modified
system. We take the depletion to be static and the exterior to provide
an ordinary reference region.

\section{From processing events to an optical metric}
\label{sec:count}

\subsection{Discrete rule and continuum assumptions}

Let $a$ be the microscopic length scale of a family of reference
networks. If a route $\gamma_a$ contains $N_a[\gamma_a]$ forwarding
events, its traversal time is
\begin{equation}
 T_a[\gamma_a]=\tau_a N_a[\gamma_a],
 \qquad \tau_a=\frac{a}{c_0}.
 \label{eq:discrete}
\end{equation}
Within each network every event has the same elementary delay $\tau_a$.
Its scaling with $a$ keeps the reference signal speed $c_0$ finite in a
continuum limit. Counting or omitting a fixed number of endpoint events
does not affect that limit.
A finite nonzero continuum advance therefore requires the number of
omitted events along a macroscopic route to be of order $1/a$; removing
only finitely many events has a vanishing effect.

For example, replacing six forwarding events by three gives
$T_a=3\tau_a$ instead of $6\tau_a$ (Fig.~\ref{fig:count}). More generally,
for corresponding routes with $N_a=(1-\ve)N_{0,a}$,
\begin{equation}
 T_a=(1-\ve)T_{0,a},
 \qquad \Delta T_a=-\ve T_{0,a}.
 \label{eq:elementary}
\end{equation}
This is the elementary time advance to be retained after coarse
graining.

\begin{figure}[t]
\centering
\setlength{\unitlength}{1mm}
\begin{picture}(82,34)
 \put(1,30){\makebox(0,0)[l]{Reference: six forwarding events}}
 \put(10,23){\line(1,0){60}}
 \multiput(10,23)(10,0){7}{\circle*{1.5}}
 \put(6,23){\makebox(0,0){$A$}}
 \put(74,23){\makebox(0,0){$B$}}
 \put(1,14){\makebox(0,0)[l]{Depleted: three forwarding events}}
 \put(10,7){\line(1,0){60}}
 \multiput(10,7)(20,0){4}{\circle*{1.5}}
 \put(6,7){\makebox(0,0){$A$}}
 \put(74,7){\makebox(0,0){$B$}}
\end{picture}
\caption{The elementary operation. Each traversal to the next retained
node incurs one delay $\tau_a$. Removing intermediate nodes and
reconnecting the sequence reduces the event count at fixed reference
endpoints. The drawing represents the microscopic rule; the continuum
model additionally requires local reconnections and isotropic
propagation.}
\label{fig:count}
\end{figure}

We make the following continuum assumptions. Reconnections are local:
their extent in reference coordinates tends to zero with $a$. On scales
large compared with $a$ but small compared with the variation scale of
the depletion, the encounter count is isotropic. For suitable microscopic
approximations $\gamma_a$ of a smooth curve $\gamma$, we assume
\begin{equation}
 aN_a[\gamma_a]\longrightarrow
 \int_\gamma n(\mathbf x)\,\dd\ell,
 \qquad \dd\ell=|\dd\mathbf x|.
 \label{eq:limit}
\end{equation}
We assume that minimum travel times converge to the minimum of this
functional as well. These are assumptions about the family of
networks, not consequences of deletion alone. For example, an
axis-aligned nearest-neighbor cubic network retains a directional
shortest-path norm when refined; refinement alone does not establish
the isotropy used in Eq.~\eqref{eq:limit}.

The dimensionless field $n$ is an \emph{encounter density along a path},
relative to the reference value. It is not, in general, the occupied
fraction of nodes per unit volume. We require $n$ to be smooth on the
continuum scale and bounded away from zero, with $n\to1$ in the
reference exterior. A depleted region has $n<1$.

\subsection{Travel times}

Equations~\eqref{eq:discrete} and~\eqref{eq:limit} give
\begin{equation}
 \begin{aligned}
 T[\gamma]&=\frac{1}{c_0}\int_\gamma n\,\dd\ell,\\
 T_n(A,B)&=\frac{1}{c_0}\inf_{\gamma:A\to B}
                  \int_\gamma n\,\dd\ell.
 \end{aligned}
 \label{eq:fermat}
\end{equation}
The optical metric is therefore
\begin{equation}
 \dd\ell_{\mathrm{opt}}^2
       =n^2(\mathbf x)\,\dd\mathbf x^2.
 \label{eq:optical}
\end{equation}
It encodes distances measured by reference travel time,
$\dd\ell_{\mathrm{opt}}=c_0\dd t$ along a ray.

There is a simple monotonicity statement. For the same endpoints and
admissible continuum paths, $n_1\leq n_2$ implies
\begin{equation}
 T_{n_1}(A,B)\leq T_{n_2}(A,B).
 \label{eq:monotonic}
\end{equation}
The inequality holds for each path integral and hence for their
infima. Thus $n\leq1$ gives a time advance or equality relative to the
reference. If the straight reference segment is admissible and passes
through $n<1$ for a nonzero length, then
$T_n(A,B)<|B-A|/c_0$.

For a weak depletion
\begin{equation}
 n(\mathbf x)=1-\ve f(\mathbf x),
 \qquad 0\leq f\leq1,\quad 0<\ve\ll1,
 \label{eq:profile}
\end{equation}
the first-order change between fixed endpoints is
\begin{equation}
 \Delta T=-\frac{\ve}{c_0}
       \int_{\gamma_0}f\,\dd\ell+O(\ve^2),
 \label{eq:timefirst}
\end{equation}
where $\gamma_0$ is the unique straight reference path. Its stationarity
makes the first-order displacement of the path irrelevant to this
expression. The remainder is understood at fixed endpoint geometry.

\section{Rays and intrinsic measurements}
\label{sec:rays}

Varying the travel-time functional, with $\mathbf e=\dd\mathbf x/\dd\ell$
the Euclidean unit tangent, gives
\begin{equation}
 \frac{\dd}{\dd\ell}(n\mathbf e)=\nabla n,
 \qquad
 \frac{\dd\mathbf e}{\dd\ell}
   =(\mathbf I-\mathbf e\mathbf e^{\mathsf T})\nabla\ln n.
 \label{eq:ray}
\end{equation}
Rays bend toward increasing encounter density. For a radially depleted
centre, $n(r)$ increases outward, so a passing ray bends away from the
centre. A weak localized depletion therefore acts as a diverging lens.
Additional central delays, for which $n$ decreases outward toward its
reference value, give the opposite deflection. Both conclusions follow
from the same processing rule.

An intrinsic observer need not describe the effect as a locally higher
signal speed. If distance is reconstructed using signal exchanges,
fewer encountered nodes mean a shorter optical distance. Observable
comparisons can use clocks and endpoints in the ordinary exterior,
or compare alternative routes with different exposure to the depleted
region. They establish relative travel times and ray geometry without
directly counting microscopic nodes. Inferring depletion as their
microscopic cause is then a model interpretation.

A uniform change that preserves all dimensionless internal
relationships can be absorbed into a choice of units. Likewise, in
one spatial dimension the optical line element becomes $\dd X^2$ under
\begin{equation}
 X(x)=\int^x n(u)\,\dd u.
 \label{eq:oned}
\end{equation}
Its variable node count alone does not establish intrinsic curvature.
In two or more spatial dimensions a position-dependent isotropic
optical metric is generally curved. Curvature depends on derivatives
of $n$, however; the inequality $n<1$ is not itself a statement that
curvature is negative everywhere.

\section{What is required for a gravitational interpretation?}
\label{sec:gravity}

\subsection{The clock prescription}

Optical data do not select a unique spacetime metric. Consider the
static family
\begin{equation}
 \dd s^2=\Omega^2(\mathbf x)
       \left[-c_0^2\dd t^2+n^2(\mathbf x)\dd\mathbf x^2\right],
 \label{eq:conformal}
\end{equation}
with $\Omega>0$ and $\Omega\to1$ in the reference exterior. Every member
has the travel-time law~\eqref{eq:fermat} and the same unparameterized
null rays. Stationary ideal clocks, in contrast, measure
\begin{equation}
 \dd\tau=\Omega\,\dd t.
 \label{eq:clock}
\end{equation}
This conformal freedom is the usual distinction between ray geometry
and clock-sensitive spacetime geometry~\cite{VisserAnalog}.

The choice $\Omega=1$ is consistent with identical clocks built from a
fixed number of equal-delay processing events. It gives the optical
effect but no clock-rate gradient; initially stationary massive
geodesics remain stationary. The choice $\Omega=1/n$ has the same
optics and position-dependent clock rates. Thus a universal clock and
probe dynamics must be supplied before optical defocusing can be
identified with repulsive free fall. Merely stipulating equal
microscopic delays does not derive that dynamics.

\subsection{A weak-field completion}

To compare the two optical signs with ordinary weak-field gravity,
adopt the static isotropic metric
\begin{equation}
 \begin{split}
 \dd s^2={}&-\left(1+\frac{2\Phi}{c_0^2}\right)c_0^2\dd t^2\\
          &+\left(1-\frac{2\Phi}{c_0^2}\right)\dd\mathbf x^2,
 \end{split}
 \label{eq:weakmetric}
\end{equation}
to first order in $|\Phi|/c_0^2$, with $\Phi\to0$ at infinity. Here the
relation between time and space factors is an additional physical
assumption. Null propagation gives
\begin{equation}
 n\simeq1-\frac{2\Phi}{c_0^2},
 \qquad
 \Phi\simeq\frac{c_0^2}{2}(1-n).
 \label{eq:potential}
\end{equation}
Accordingly, added delays give a negative potential, whereas depletion
gives a positive one relative to the exterior.

If clocks and massive probes couple to this same metric, stationary
clocks obey $\dd\tau/\dd t\simeq1+\Phi/c_0^2$, and slow freely falling
probes obey $\mathbf a=-\nabla\Phi$. For Eq.~\eqref{eq:profile},
\begin{equation}
 \Phi=\frac{c_0^2\ve}{2}f,
 \qquad
 \mathbf a=-\frac{c_0^2\ve}{2}\nabla f
 \label{eq:acceleration}
\end{equation}
at leading order. A central depletion with $f'(r)<0$ is a potential
hill and gives $a_r>0$: outward free fall. Clocks there run faster
relative to exterior clocks. Locally, ideal rods and clocks defined by
the metric still measure the characteristic signal speed $c_0$.

The factor $1/2$ in Eq.~\eqref{eq:potential} comes from the particular
time--space relation in Eq.~\eqref{eq:weakmetric}; it is not fixed by
node counting. A different clock prescription can retain the same
optical index and change the gravitational interpretation. Equations
\eqref{eq:weakmetric}--\eqref{eq:acceleration} specify a conditional
completion of the model, not a derivation of Einstein's equations or
of an equivalence principle.

\section{A smooth example}
\label{sec:example}

Consider a spherically symmetric encounter deficit
\begin{equation}
 n(r)=1-\ve e^{-r^2/L^2},
 \qquad 0<\ve\ll1.
 \label{eq:gauss}
\end{equation}
It is regular at the origin and approaches the reference exponentially.
The following observables illustrate the same depletion in two
different optical experiments.

First, place endpoints at $(\mathbf b,-Z)$ and $(\mathbf b,Z)$, with
$b=|\mathbf b|$ and $Z$ fixed. The reference path is parallel to the
$z$ axis. Equation~\eqref{eq:timefirst} gives
\begin{equation}
 \Delta T=-\frac{\sqrt{\pi}\ve L}{c_0}
       e^{-b^2/L^2}\operatorname{erf}(Z/L)+O(\ve^2).
 \label{eq:gausstime}
\end{equation}
The coefficient of the advance approaches a finite limit for distant
endpoints. This is a fixed-endpoint perturbation calculation; it does
not compare arrival times of rays ending in different directions at
infinity.

Second, send in a ray from the ordinary exterior with impact vector
$\mathbf b$. Integrating Eq.~\eqref{eq:ray} along its unperturbed
trajectory gives the asymptotic change in direction
\begin{equation}
 \Delta\mathbf e_\perp
  =2\sqrt{\pi}\ve\frac{\mathbf b}{L}
       e^{-b^2/L^2}+O(\ve^2).
 \label{eq:gaussbend}
\end{equation}
The deflection points away from the centre. A central ray has zero
deflection by symmetry, and rays far outside the depleted region are
unaffected to exponential accuracy.

Under the gravitational completion of Sec.~\ref{sec:gravity},
\begin{equation}
 \Phi(r)=\frac{c_0^2\ve}{2}e^{-r^2/L^2},
 \qquad
 a_r=\frac{c_0^2\ve r}{L^2}e^{-r^2/L^2}.
 \label{eq:gaussforce}
\end{equation}
These expressions are again leading-order relations. They give a
regular outward acceleration for every $r>0$.

One may additionally ask what source an observer using the Newtonian
field equation would infer. Defining that source by
\begin{equation}
 \nabla^2\Phi=4\pi G\reff
 \label{eq:poisson}
\end{equation}
gives, in three spatial dimensions,
\begin{equation}
 \reff(r)=\frac{c_0^2\ve}{8\pi G L^2}
      \left(\frac{4r^2}{L^2}-6\right)e^{-r^2/L^2}.
 \label{eq:gaussrho}
\end{equation}
The inferred source is negative in the centre and positive in an outer
region. Its total integral vanishes at the order considered. The enclosed inferred
mass is
\begin{equation}
 \Meff(r)=\frac{r^2\Phi'(r)}{G}
       =-\frac{c_0^2\ve r^3}{G L^2}e^{-r^2/L^2},
 \label{eq:enclosed}
\end{equation}
which is negative for finite $r>0$ and tends to zero at infinity.
Thus a local depleted profile can give outward acceleration without
representing an isolated negative total mass. The fitted density
$\reff$ is defined by Eq.~\eqref{eq:poisson}; it is not the microscopic
number density of removed nodes or a calculation of their energy.

\section{Long-range behavior and the remaining physical question}
\label{sec:discussion}

A negative-mass type exterior has a more specific requirement. Write
$\eta=1-n$. If, outside a core,
\begin{equation}
 \eta(r)\simeq\frac{\ell_*}{r},\qquad \ell_*>0,
 \label{eq:tail}
\end{equation}
then the weak-field completion gives
\begin{equation}
 \Phi(r)\simeq\frac{c_0^2\ell_*}{2r},
 \qquad
 a_r\simeq\frac{c_0^2\ell_*}{2r^2}.
 \label{eq:inverse}
\end{equation}
This has the exterior form of a negative effective mass
\begin{equation}
 \Meff=-\frac{c_0^2\ell_*}{2G},
 \label{eq:negative}
\end{equation}
in the regime $r\gg\ell_*$. An exterior with exactly $n=1$ has no such
tail. For a smooth compactly supported deficit, the inferred total
source in Eq.~\eqref{eq:poisson} is zero by the divergence theorem.

The long-range profile is additional physics. A finite set of removed
nodes confined to a bounded region does not by itself establish
Eq.~\eqref{eq:tail}. It would have to produce an extended change in the
encounter density or in the effective processing rule of the remaining
substrate. A microscopic theory would need to determine that response,
its stability, and the dynamics of the objects used as clocks and
probes. It would also have to establish that relevant signal species
share the same continuum geometry. Frequency-dependent or
species-dependent delays would give further observable structure.

The result established by the counting rule is narrower and direct:
local contraction of processing sequences gives an optical index below
the reference value, earlier traversal, and outward deflection for a
central depleted profile. Applying the usual weak-field metric
relation turns the same profile into faster interior clocks and
repulsive free fall. The first set of statements follows from the
continuum propagation model; the second identifies a precise additional
hypothesis under which the missing-step intuition becomes a
gravitational analogue.

\begin{thebibliography}{3}

\bibitem{kleinert-2004}
H. Kleinert and J. Zaanen,
``Nematic world crystal model of gravity explaining absence of torsion
in spacetime,''
\href{https://doi.org/10.1016/j.physleta.2004.03.048}
{Phys. Lett. A \textbf{324}, 361--365 (2004)}.

\bibitem{toffoli:79}
T. Toffoli, ``The Role of the Observer in Uniform Systems,''
in \emph{Applied General Systems Research: Recent Developments and
Trends}, edited by G. J. Klir (Plenum Press, New York, 1978),
pp.~395--400,
\href{https://doi.org/10.1007/978-1-4757-0555-3_29}
{doi:10.1007/978-1-4757-0555-3\_29}.

\bibitem{VisserAnalog}
M. Visser, ``Acoustic black holes: horizons, ergospheres and Hawking
radiation,''
\href{https://doi.org/10.1088/0264-9381/15/6/024}
{Class. Quantum Grav. \textbf{15}, 1767--1791 (1998)}.

\end{thebibliography}
\end{document}
